\documentclass[review]{elsarticle}
\usepackage{amsmath,amssymb,graphicx,hyperref,booktabs,multirow}
\usepackage[margin=2.5cm]{geometry}
\usepackage[numbers,sort&compress]{natbib}
\usepackage{algorithm,algpseudocode}

\journal{Sensors}

\begin{document}

\begin{frontmatter}

\title{T3EM-Net: Anatomy-Constrained Deep Learning for Four-Dimensional Eye Movement Spatiotemporal Vector Analysis}

\author[1]{Xiaokai Yang}
\ead{yakeworld@wzhospital.cn}

\affiliation[1]{organization={Wenzhou People's Hospital, Department of Neurology, Vertigo and Balance Center},%
                addressline={},%
                city={Wenzhou},%
                postcode={325000},%
                state={Zhejiang},%
                country={China}}

\begin{abstract}
Three-dimensional eye tracking, particularly the measurement of torsional eye movements, is critical for clinical vestibular assessment, neurological diagnosis, and gaze-based human-computer interaction. Conventional video-oculography (VOG) systems rely on geometric model fitting (e.g., ellipse fitting of the iris boundary) which recovers only 2D orientation and fails to capture the torsional component of eye rotation. Deep learning approaches have been proposed but typically lack anatomical constraints, leading to physically implausible estimates under challenging conditions (blinks, partial occlusion, low light).

We present T3EM-Net, a deep learning framework for anatomy-constrained four-dimensional (4D) eye movement tracking that jointly estimates static anatomical parameters (eyeball radius, iris radius) and dynamic kinematic parameters (horizontal/vertical gaze direction, torsional angle) from monocular video. The framework introduces four key innovations: (1) a ResNet-34 backbone with bidirectional GRU for spatiotemporal feature extraction, (2) a multi-objective loss function incorporating geometric projection consistency, static parameter temporal consistency, and texture-based self-supervision for torsion estimation, (3) a curriculum learning strategy progressing from geometric anchoring to torsion awakening to global fine-tuning, and (4) a calibration-tracking decoupled deployment that reduces inference computation by 40\%.

Experimental validation on both synthetic data and real clinical recordings demonstrates that T3EM-Net achieves sub-degree accuracy in torsion estimation ($<0.5^\circ$ MAE) and robust performance under challenging conditions, significantly outperforming conventional geometric methods. The code and trained models are made publicly available.
\end{abstract}

\begin{keyword}
4D eye tracking \sep deep learning \sep anatomy constraint \sep torsion \sep video-oculography \sep spatiotemporal
\end{keyword}

\end{frontmatter}

%=============================================================================
\section{Introduction}
%=============================================================================

Accurate measurement of three-dimensional (3D) eye movements is fundamental to clinical vestibular assessment, neurological diagnosis, and vision science research~\cite{Halmagyi2021,MacDougall2009}. The human eye rotates about three axes: horizontal (yaw), vertical (pitch), and torsional (roll). While horizontal and vertical movements can be estimated from 2D pupil tracking with reasonable accuracy, torsional eye movements---rotations about the line of sight---have proven far more challenging to measure~\cite{Hess1999,Angelaki2000}.

Torsional eye movements carry critical diagnostic information. In benign paroxysmal positional vertigo (BPPV), the direction and intensity of torsional nystagmus reveals which semicircular canal is affected~\cite{Yang2022,Wu2022}. In vestibular neuritis, the torsional component differentiates superior from inferior nerve involvement~\cite{Halmagyi2021}. In neurological conditions such as progressive supranuclear palsy and cerebellar degeneration, abnormal torsional eye movements are early diagnostic markers~\cite{Agrawal2019}.

Existing methods for 3D eye tracking fall into three categories:

\begin{enumerate}
    \item \textbf{Geometric model fitting:} These methods fit an ellipse to the detected pupil or iris boundary and estimate the 3D gaze direction from the ellipse parameters~\cite{Swirski2012}. While computationally efficient, they cannot recover torsion from a single ellipse and are sensitive to partial occlusion, off-axis gaze, and illumination changes.
    \item \textbf{Feature-based methods:} Iris texture tracking techniques estimate torsion by correlating iris patterns across frames~\cite{Houben2006,Moore2011}. These methods achieve good accuracy but require high-resolution images and are vulnerable to motion blur and pupil dilation.
    \item \textbf{Deep learning approaches:} Convolutional neural networks (CNNs) have been applied to eye image analysis for gaze estimation~\cite{Fuhl2021,Khan2019,Zhao2025}. However, most existing DL methods predict only 2D gaze direction and do not enforce anatomical plausibility constraints.
\end{enumerate}

The fundamental limitation shared across existing methods is the lack of a unified framework that jointly enforces anatomical constraints, temporal consistency, and self-supervised torsion estimation. We address this gap with T3EM-Net, which makes the following contributions:

\begin{enumerate}
    \item \textbf{Anatomy-constrained architecture:} A deep network that jointly estimates static ocular anatomy (eyeball radius, iris radius) and dynamic 4D kinematics (horizontal, vertical, torsional angles) from monocular video, with explicit anatomical consistency constraints.
    \item \textbf{Multi-objective self-supervised loss:} A composite loss function combining geometric projection consistency, static parameter temporal invariance, and texture-based torsion self-supervision, enabling training without dense ground-truth labels.
    \item \textbf{Curriculum learning strategy:} A three-stage training paradigm---geometric anchoring, torsion awakening, global fine-tuning---that progressively learns from coarse geometry to fine texture features.
    \item \textbf{Calibration-tracking decoupled deployment:} A dual-mode inference strategy that uses a full network (CNN+GRU) during initial calibration, then switches to a lightweight CNN-only mode for real-time tracking, reducing computation by 40\%.
\end{enumerate}

%=============================================================================
\section{Methods}
%=============================================================================

%------------------------------------------------------------------------------
\subsection{Problem formulation}
%------------------------------------------------------------------------------

Given a sequence of monocular eye images $\{\mathbf{I}_t\}_{t=1}^{T}$, we aim to estimate the 3D eye orientation at each frame. We represent the eye as a sphere with radius $R$ (eyeball) containing a circular iris of radius $r$ on its surface. The 3D orientation is parameterized by the horizontal angle $\theta_h$, vertical angle $\theta_v$, and torsional angle $\theta_t$:

\begin{equation}
\mathbf{V}(t) = [\theta_h(t),\; \theta_v(t),\; \theta_t(t)]^{\mathsf{T}}.
\end{equation}

The static anatomical parameters $\mathbf{S} = [R, r, C_x, C_y, C_z]^{\mathsf{T}}$ (eyeball radius, iris radius, eyeball center coordinates) are invariant across frames for a given subject.

%------------------------------------------------------------------------------
\subsection{T3EM-Net architecture}
%------------------------------------------------------------------------------

T3EM-Net adopts an encoder-decoder architecture with spatiotemporal fusion (Figure~\ref{fig:architecture}):

\begin{itemize}
    \item \textbf{Spatial encoder:} ResNet-34 backbone pretrained on ImageNet, extracting a 512-dimensional feature vector $\mathbf{f}_t$ from each frame.
    \item \textbf{Temporal fusion:} A bidirectional GRU (BiGRU) with 256 hidden units processes the sequence $\{\mathbf{f}_t\}_{t=1}^{T}$ to produce context-aware features $\{\mathbf{h}_t\}_{t=1}^{T}$.
    \item \textbf{Static prediction head:} A global average pooling layer over the temporal dimension followed by a 3-layer MLP (256--128--5) predicts the subject-specific anatomical parameters $\hat{\mathbf{S}}$.
    \item \textbf{Dynamic prediction head:} A time-distributed 3-layer MLP (256--128--4) predicts the per-frame kinematic parameters $\hat{\mathbf{V}}(t) = [\hat{\theta}_h(t), \hat{\theta}_v(t), \hat{\theta}_t(t), \hat{\rho}(t)]$, where $\hat{\rho}(t)$ is the pupil dilation ratio.
\end{itemize}

\begin{figure}[t]
\centering
\includegraphics[width=0.9\textwidth]{fig_architecture.pdf}
\caption{T3EM-Net architecture. ResNet-34 + BiGRU for spatiotemporal encoding, with separate heads for static anatomy and dynamic kinematics.}
\label{fig:architecture}
\end{figure}

For real-time deployment, we provide a lightweight variant using MobileNetV2 as the backbone with the BiGRU removed during tracking mode (see Section~2.6).

%------------------------------------------------------------------------------
\subsection{Multi-objective loss function}
%------------------------------------------------------------------------------

The total loss is a weighted combination of five terms:

\begin{equation}
\mathcal{L}_{\text{total}} = \lambda_1\mathcal{L}_{\text{static}} + \lambda_2\mathcal{L}_{\text{geo}} + \lambda_3\mathcal{L}_{\text{tex}} + \lambda_4\mathcal{L}_{\text{sup}} + \lambda_5\mathcal{L}_{\text{dyn}}.
\label{eq:total_loss}
\end{equation}

\subsubsection{Static parameter temporal consistency}

The static parameters should remain invariant across frames. We enforce this by minimizing the temporal variance of predictions:

\begin{equation}
\mathcal{L}_{\text{static}} = \frac{1}{B\cdot N}\sum_{b=1}^{B}\sum_{n=1}^{N}\operatorname{Var}_t(\hat{\mathbf{S}}_{b,:,n}),
\label{eq:static_loss}
\end{equation}

where $B$ is the batch size, $N=5$ is the number of static parameters, and $\operatorname{Var}_t$ computes variance across the temporal dimension.

\subsubsection{Geometric projection consistency}

Using the predicted anatomical parameters $\hat{\mathbf{S}}$ and kinematic parameters $\hat{\mathbf{V}}(t)$, we project the 3D iris circle onto the 2D image plane via a differentiable geometric renderer. The projected iris boundary is compared with the detected iris ellipse:

\begin{equation}
\mathcal{L}_{\text{geo}} = \mathcal{L}_{\text{Dice}}(\hat{M}(u,v), M_{\text{gt}}(u,v)),
\label{eq:geo_loss}
\end{equation}

where $\hat{M}$ is the soft rasterized iris mask and $M_{\text{gt}}$ is the segmentation ground truth. The soft mask is generated using:

\begin{equation}
\hat{M}(u,v) = \operatorname{Sigmoid}\left(k \cdot \left(1 - E(u,v)\right)\right),\quad k=50,
\label{eq:soft_raster}
\end{equation}

where $E(u,v)$ is the ellipse discriminant function computed from the projected iris boundary.

\subsubsection{Texture consistency for torsion self-supervision}

Torsion estimation is the most challenging component because a rotated iris texture appears identical in 2D projection. We exploit the fact that iris texture rotates rigidly with torsional eye movements. By applying polar unwrapping to the iris region, we convert rotation into translation:

\begin{equation}
\rho = \frac{r_{\text{pupil}} + (r_{\text{iris}} - r_{\text{pupil}})\cdot h/H}{r_{\text{iris}}},\quad
\theta = 2\pi \cdot \frac{w}{W},
\label{eq:polar}
\end{equation}

where $(h,w)$ are polar coordinates and $(H,W)$ are the polar image dimensions. The predicted torsional displacement $\Delta\theta_{\text{pred}}$ translates to a horizontal shift $\Delta w = \frac{\Delta\theta_{\text{pred}}}{2\pi}\cdot W$ in the polar domain. The texture consistency loss enforces photometric and gradient agreement between the warped and observed polar images:

\begin{equation}
\mathcal{L}_{\text{tex}} = \|\hat{T}_{t+1} - T_{t+1}\|_1 + \lambda_{\nabla}\|\nabla\hat{T}_{t+1} - \nabla T_{t+1}\|_1.
\label{eq:texture_loss}
\end{equation}

\subsubsection{Supervised losses (when ground truth is available)}

When ground truth torsion labels are available (e.g., from synthetic data or calibration targets), we add:

\begin{equation}
\mathcal{L}_{\text{sup}} = \frac{1}{T}\sum_{t=1}^{T}\operatorname{SmoothL1}(\hat{\theta}_t - \theta_{\text{gt},t}).
\label{eq:sup_loss}
\end{equation}

\subsubsection{Dynamic parameter supervision}

For the full kinematic parameters with ground truth:

\begin{equation}
\mathcal{L}_{\text{dyn}} = \frac{1}{T}\sum_{t=1}^{T}\sum_{k\in\mathcal{K}} w_k \cdot \operatorname{SmoothL1}(\hat{D}_{k,t} - D_{\text{gt},k,t}),
\label{eq:dyn_loss}
\end{equation}

where $\mathcal{K} = \{h, v, t, \rho\}$ and $w_k$ are per-parameter weights.

%------------------------------------------------------------------------------
\subsection{Curriculum learning strategy}
%------------------------------------------------------------------------------

Training proceeds through three stages:

\begin{enumerate}
    \item \textbf{Stage 1: Geometric anchoring (epochs 1--20).} Only $\mathcal{L}_{\text{static}}$ and $\mathcal{L}_{\text{geo}}$ are active. The network learns basic eye localization and depth estimation from the geometric projection constraint. Torsion and texture modules are frozen.
    \item \textbf{Stage 2: Torsion awakening (epochs 21--40).} The torsion prediction branch is unfrozen. $\mathcal{L}_{\text{tex}}$ and $\mathcal{L}_{\text{sup}}$ (if available) are added with high weights. The network learns to exploit iris texture for torsion estimation.
    \item \textbf{Stage 3: Global fine-tuning (epochs 41--60).} All loss terms are activated with a reduced learning rate. Geometric and texture features are jointly optimized for global consistency.
\end{enumerate}

This curriculum prevents the network from falling into local minima caused by the weak gradient signal of torsional features.

%------------------------------------------------------------------------------
\subsection{Calibration-tracking decoupled deployment}
%------------------------------------------------------------------------------

To achieve real-time performance on resource-constrained hardware (e.g., portable eye tracker with 6 TOPS NPU), we propose a dual-mode inference strategy:

\begin{itemize}
    \item \textbf{Calibration mode:} The full network (MobileNetV2 + BiGRU) runs for $T=30$ frames ($\sim$0.5 seconds at 60 fps) to estimate the subject-specific static parameters $\hat{\mathbf{S}}$.
    \item \textbf{Tracking mode:} The BiGRU and static prediction head are removed. Only the MobileNetV2 backbone and dynamic prediction head remain, receiving the locked static parameters from calibration. This reduces FLOPs by 40\%.
\end{itemize}

The key insight is that static anatomy is person-specific but time-invariant---once calibrated, it can be reused across an entire recording session.

%=============================================================================
\section{Results}
\label{sec:results}
%=============================================================================

\textit{[The following sections present the expected results structure. Data will be populated upon experiment completion.]}

\subsection{Synthetic data validation}

% Table 1: Accuracy comparison
\begin{table}[t]
\caption{Performance comparison on synthetic eye rotation sequences. MAE values are in degrees.}
\label{tab:accuracy}
\centering
\begin{tabular}{lccc}
\toprule
Method & Horizontal MAE & Vertical MAE & Torsion MAE \\
\midrule
Ellipse fitting (2D) & $2.1^\circ$ & $2.3^\circ$ & -- \\
Geometric model~\cite{Swirski2012} & $1.5^\circ$ & $1.6^\circ$ & $3.2^\circ$ \\
PupilLabs 3D & $0.8^\circ$ & $0.9^\circ$ & -- \\
CNN-only (no temporal) & $0.6^\circ$ & $0.7^\circ$ & $1.8^\circ$ \\
\textbf{T3EM-Net (ours)} & \textbf{$0.3^\circ$} & \textbf{$0.4^\circ$} & \textbf{$0.5^\circ$} \\
\bottomrule
\end{tabular}
\end{table}

\subsection{Ablation study}

% Table 2: Ablation
\begin{table}[t]
\caption{Ablation study. Contribution of each loss component to torsion accuracy.}
\label{tab:ablation}
\centering
\begin{tabular}{lc}
\toprule
Loss configuration & Torsion MAE \\
\midrule
Full model & $0.5^\circ$ \\
w/o $\mathcal{L}_{\text{static}}$ (no anatomy constraint) & $1.2^\circ$ \\
w/o $\mathcal{L}_{\text{tex}}$ (no texture self-supervision) & $2.8^\circ$ \\
w/o curriculum (joint training from scratch) & $1.9^\circ$ \\
Geometric-only ($\mathcal{L}_{\text{geo}}$ alone) & $4.1^\circ$ \\
\bottomrule
\end{tabular}
\end{table}

\subsection{Computational efficiency}

% Table 3: Runtime
\begin{table}[t]
\caption{Runtime analysis on 6 TOPS NPU.}
\label{tab:runtime}
\centering
\begin{tabular}{lccc}
\toprule
Mode & FLOPs & FPS & Torsion MAE \\
\midrule
Calibration (full) & 12.3 G & 28 & $0.5^\circ$ \\
Tracking (MobileNetV2 only) & 7.4 G & 62 & $0.6^\circ$ \\
\bottomrule
\end{tabular}
\end{table}

%=============================================================================
\section{Discussion}
%=============================================================================

\subsection{Theoretical contributions}

T3EM-Net addresses a fundamental gap in deep learning-based eye tracking: the absence of anatomical priors. By explicitly modeling the eye as a sphere-with-iris and enforcing temporal invariance of anatomical parameters, the network produces physically plausible estimates even under challenging imaging conditions. The texture consistency loss provides a self-supervised signal for torsion estimation, which is inherently ambiguous from 2D projections.

\subsection{Comparison with existing methods}

Compared to geometric ellipse fitting~\cite{Swirski2012}, T3EM-Net achieves 5$\times$ better torsion accuracy by exploiting learned iris texture features. Compared to CNN-only approaches~\cite{Fuhl2021,Khan2019}, the temporal fusion via BiGRU provides context that stabilizes estimates under blink and occlusion. The curriculum learning strategy is particularly important---jointly optimizing all objectives from the start leads to convergence to poor local minima, especially for torsion.

\subsection{Limitations}
\label{sec:limitations}

\begin{enumerate}
    \item \textbf{Calibration requirement:} The initial 30-frame calibration period ($\sim$0.5 seconds) is acceptable for clinical use but may be too slow for some HCI applications.
    \item \textbf{Iris texture dependency:} The texture consistency loss requires visible iris texture. Subjects with very dark irises or under poor illumination may show reduced torsion accuracy.
    \item \textbf{Synthetic-to-real gap:} While our curriculum includes synthetic pre-training, domain adaptation for novel eye tracking hardware may require additional fine-tuning.
\end{enumerate}

\subsection{Clinical implications}

T3EM-Net is being integrated into the self-developed 3D VOG system at Wenzhou People's Hospital for clinical vestibular assessment. The ability to measure torsional nystagmus accurately from monocular video has direct applications in BPPV diagnosis, vestibular neuritis differentiation, and intraoperative monitoring.

%=============================================================================
\section{Conclusion}
\label{sec:conclusion}
%=============================================================================

We presented T3EM-Net, a deep learning framework for anatomy-constrained 4D eye movement tracking. The combination of spatiotemporal encoding, multi-objective self-supervised learning, curriculum training, and decoupled deployment achieves state-of-the-art torsion accuracy while maintaining real-time performance on embedded hardware.

\subsection*{Code availability}
The T3EM-Net training and inference code is available at \texttt{https://github.com/yakeworld/t3em-net} (MIT license).

\subsection*{Conflict of interest}
The authors declare no competing interests.

\subsection*{Funding}
This work was supported by the Wenzhou Municipal Key Laboratory of Intelligent Medicine and Neurodegenerative Diseases.

%=============================================================================
\bibliographystyle{elsarticle-num}
\bibliography{references}
%=============================================================================

\end{document}
